\section{Introduction}\label{sec:intro}

Two of the defining challenges of this century compete for the same resource: land. The energy transition requires terawatts of new solar photovoltaic (PV) capacity, and ground-mounted solar farms need roughly one to two hectares per megawatt. Food production must rise while crops are exposed to more frequent heat waves and droughts. Warming alone is projected to reduce global yields of maize, wheat, rice and soybean by several per cent per degree \citep{zhao2017}, and a new generation of climate and crop models shows these impacts emerging earlier than previously expected \citep{jagermeyr2021}. Croplands are also where solar power potential is greatest \citep{adeh2019}. They are flat, sunny, accessible and close to rural electricity demand, so converting farmland into solar farms is tempting and controversial at the same time.

Agrivoltaic (AV) systems combine crops and PV modules on the same land and aim to resolve this conflict \citep{dupraz2011,dinesh2016,weselek2019}. By sharing sunlight between modules and plants, a hectare can deliver more food plus energy than the two uses side by side, with land-equivalent ratios (LER) well above one in field trials and modelling studies \citep{dupraz2011,amaducci2018,valle2017,trommsdorff2021}. The partial shade also changes the crop microclimate. It lowers canopy temperature and evaporative demand, conserves soil moisture and improves water-use efficiency, especially in drylands \citep{marrou2013,adeh2018,barron2019,elamri2018,weselek2021}. Agrivoltaics is therefore a candidate adaptation measure against heat and drought stress, in addition to a mitigation measure through renewable electricity. The same shade, however, also reduces the light available for photosynthesis. Regulators increasingly require AV systems to retain most of the reference crop yield. The German pre-standard DIN SPEC 91434, for example, requires at least 66\% of the reference yield \citep{din2021}, and other schemes and many farmers set stricter floors. How much electricity can be produced without sacrificing food security depends on the crop, the local climate, the year and the growth stage.

Most AV designs fix the geometry of the array, i.e.\ the ground coverage ratio (GCR, module area per ground area) and the tilt or tracking rule, from generic rules of thumb. Single-axis trackers add a degree of freedom: the modules can be rotated away from the sun to let light through when the crop needs it, as demonstrated for lettuce by \citet{valle2017}, or towards the sun to maximise electricity and shade the crop during heat waves. Choosing the rotation every day is a sequential decision problem. The value of light depends on the crop's phenological stage, soil water and upcoming weather, and the effect of a decision appears in the harvest months later. Designing the GCR and the control policy for every combination of crop and climate on Earth, and for a warming climate, is beyond manual rule design. Recent work in this journal has used evolutionary multi-objective optimisation to lay out PV arrays on farmland \citep{zhang2026}. Machine learning for crop management has mainly used reinforcement learning with black-box simulators \citep{gautron2022}, which is sample-inefficient and hard to scale to global, multi-decadal assessments.

This paper proposes \emph{SALS} (stress-aware light sharing), a learning algorithm that designs and controls agrivoltaic systems jointly, for any cropland site on Earth, through a differentiable physics simulator. The simulator couples hourly solar geometry, single-axis trackers with backtracking, row-to-row shading and PV electricity, the direct and diffuse light reaching the crop, a microclimate model of canopy temperature and reference evapotranspiration under the modules, and the SIMPLE crop model with heat and drought stress \citep{zhao2019}. Because every component is written as differentiable tensor operations, two small neural networks can be trained by gradient descent through entire growing seasons. A \emph{design network} maps site descriptors to the GCR, and a \emph{control network} maps the daily crop and soil state and the same-day weather forecast to the tracker's light-sharing level. Training maximises electricity subject to a yield-retention floor, across hundreds of sites, twenty years of weather and four warming levels. This physics-informed learning approach \citep{karniadakis2021} gives a single policy that generalises to unseen regions, years and climates.

We apply SALS to 400 sites that represent where maize, wheat, rice, soybean and potato are grown worldwide, sampled in proportion to their harvested area \citep{yu2020} with local crop calendars \citep{jagermeyr2021}, and driven by NASA POWER daily weather for 2001--2019 \citep{sparks2018}. The contributions are as follows.
\begin{enumerate}
\item A differentiable agrivoltaic simulator that couples tracker geometry, PV output, crop-level light and microclimate, soil water and a generic crop model, and runs thousands of site-seasons in parallel on a CPU.
\item The SALS algorithm, which learns the array design and a closed-loop, stress-aware tracker policy end-to-end under a food-security constraint. It is evaluated on geographically held-out sites and years against static designs, expert rules and a perfect-foresight oracle.
\item A global assessment under warming of +1.5, +2 and +3\,\textdegree C, which quantifies how much electricity croplands can supply while retaining at least 90\% of crop yields, and where agrivoltaics protects yields against heat and drought.
\end{enumerate}
